\documentclass[11pt]{article}

\usepackage[spanish]{babel}
\usepackage[utf8]{inputenc}
\usepackage[T1]{fontenc}
\usepackage{amsmath, amssymb}
\usepackage{geometry}
\usepackage{array}
\usepackage{booktabs}
\usepackage{multicol}

\geometry{margin=1.5cm}

\newcommand{\jw}{j\omega}
\newcommand{\degr}{^\circ}

\begin{document}

\begin{center}
{\Large \textbf{Formulario: Respuesta en Frecuencia y Bode}}\\
\vspace{0.2cm}
\textbf{Magnitud, fase, pendientes y respuesta senoidal}
\end{center}

\begin{multicols}{2}

% ============================
\section*{1. Sustitución en frecuencia}

Para analizar una función de transferencia en frecuencia:
\[
s = j\omega
\]
Entonces:
\[
G(s) \rightarrow G(j\omega)
\]
Ejemplo:
\[
G(s)=\frac{1}{s}
\]
\[
G(j\omega)=\frac{1}{j\omega}
\]

% ============================
\section*{2. Números complejos}

Un número complejo tiene forma:
\[
z=a+jb
\]

\subsection*{Magnitud}
\[
|z|=\sqrt{a^2+b^2}
\]
Ejemplos:
\[
|1+j\omega|=\sqrt{1+\omega^2}
\]
\[
|5+j\omega|=\sqrt{25+\omega^2}
\]
\[
|j\omega|=\omega
\]

\subsection*{Fase}
\[
\angle z=\tan^{-1}\!\left(\frac{b}{a}\right)
\]
Ejemplos:
\[
\angle(1+j\omega)=\tan^{-1}(\omega)
\]
\[
\angle(5+j\omega)=\tan^{-1}\!\left(\frac{\omega}{5}\right)
\]
\[
\angle(j\omega)=90^\circ
\]
\[
\angle(-j\omega)=-90^\circ
\]

\subsection*{Conjugado}

Si $z=a+jb$, su conjugado es $\overline{z}=a-jb$.

Para dividir complejos:
\[
\frac{a+jb}{c+jd}\cdot\frac{c-jd}{c-jd}
\]
Propiedad importante:
\[
(c+jd)(c-jd)=c^2+d^2
\]

% ============================
\section*{3. Magnitud de una función}

Si:
\[
G(j\omega)=\frac{A(j\omega)}{B(j\omega)}
\]
entonces:
\[
|G(j\omega)|=\frac{|A(j\omega)|}{|B(j\omega)|}
\]
Si hay multiplicaciones:
\[
|ABC|=|A||B||C|
\]
Ejemplo:
\[
G(s)=\frac{2(s+1)}{s(s+5)}
\]
\[
G(j\omega)=\frac{2(1+j\omega)}{j\omega(5+j\omega)}
\]
\[
|G(j\omega)|=\frac{2|1+j\omega|}{|j\omega||5+j\omega|}
\]
\[
|G(j\omega)|=\frac{2\sqrt{1+\omega^2}}{\omega\sqrt{25+\omega^2}}
\]

% ============================
\section*{4. Fase de una función}

Reglas:
\[
\angle(AB)=\angle A+\angle B
\]
\[
\angle\!\left(\frac{A}{B}\right)=\angle A-\angle B
\]
Si:
\[
G(j\omega)=\frac{2(1+j\omega)}{j\omega(5+j\omega)}
\]
entonces:
\[
\angle G(j\omega)=\angle 2+\angle(1+j\omega)-\angle(j\omega)-\angle(5+j\omega)
\]
Como $\angle 2=0^\circ$, $\angle(j\omega)=90^\circ$:
\[
\angle G(j\omega)=\tan^{-1}(\omega)-90^\circ-\tan^{-1}\!\left(\frac{\omega}{5}\right)
\]

% ============================
\section*{5. Decibeles}

La magnitud en decibeles es:
\[
|G(j\omega)|_{\mathrm{dB}}=20\log_{10}|G(j\omega)|
\]
Reglas útiles:
\[
20\log_{10}(AB)=20\log_{10}(A)+20\log_{10}(B)
\]
\[
20\log_{10}\!\left(\frac{A}{B}\right)=20\log_{10}(A)-20\log_{10}(B)
\]
\[
20\log_{10}(10)=20 \qquad 20\log_{10}(0.1)=-20
\]

% ============================
\section*{6. Forma estándar de Bode}

La forma estándar busca escribir factores así:
\[
1+\frac{s}{a}
\]
Ejemplos:
\[
s+10=10\!\left(1+\frac{s}{10}\right), \quad s+2=2\!\left(1+\frac{s}{2}\right)
\]
\[
s+25=25\!\left(1+\frac{s}{25}\right), \quad s+40=40\!\left(1+\frac{s}{40}\right)
\]
La función queda como:
\[
G(s)=K\,s^n\,\frac{\text{ceros normalizados}}{\text{polos normalizados}}
\]
La ganancia pura es:
\[
K=\frac{\text{constantes del numerador}}{\text{constantes del denominador}}
\]
Ejemplo:
\[
G(s)=\frac{20s(s+10)}{(s+2)(s+25)(s+40)}
\]
\[
K=\frac{20(10)}{2(25)(40)}=\frac{200}{2000}=0.1
\]

% ============================
\section*{7. Frecuencias de quiebre}

Un factor $s+a$ o $1+s/a$ produce una frecuencia de quiebre en $\omega=a$.

Ejemplos:
\[
s+2 \Rightarrow \omega=2, \quad s+10 \Rightarrow \omega=10
\]
\[
s+25 \Rightarrow \omega=25, \quad s+40 \Rightarrow \omega=40
\]
Si aparece $(s+1)^2$, la frecuencia de quiebre es $\omega=1$ pero cuenta como dos polos o dos ceros.

% ============================
\section*{8. Pendientes en Bode de magnitud}

\begin{center}
\begin{tabular}{lc}
\toprule
Elemento & Cambio de pendiente \\
\midrule
Cero & $+20$ dB/dec \\
Polo & $-20$ dB/dec \\
Cero en el origen & $+20$ dB/dec desde el inicio \\
Polo en el origen & $-20$ dB/dec desde el inicio \\
Cero doble & $+40$ dB/dec \\
Polo doble & $-40$ dB/dec \\
\bottomrule
\end{tabular}
\end{center}

Regla rápida:
\[
s^n \text{ arriba} \Rightarrow +20n\ \text{dB/dec}
\]
\[
s^n \text{ abajo} \Rightarrow -20n\ \text{dB/dec}
\]

% ============================
\section*{9. Actualizar anclajes de magnitud}

Si ya conoces la magnitud en $\omega_1$, puedes calcularla en $\omega_2$ con:
\[
|G(j\omega_2)|_{\mathrm{dB}}=|G(j\omega_1)|_{\mathrm{dB}}+m\log_{10}\!\left(\frac{\omega_2}{\omega_1}\right)
\]
donde $m$ es la pendiente en dB/dec.

Ejemplo: $|G(j40)|=-6\ \text{dB}$, pendiente $m=-20$, para $\omega=100$:
\[
|G(j100)|_{\mathrm{dB}}=-6-20\log_{10}\!\left(\frac{100}{40}\right)=-6-20\log_{10}(2.5)\approx -14\ \text{dB}
\]

% ============================
\section*{10. Anclaje inicial}

Si no hay polos ni ceros en el origen:
\[
|G(j\omega)|_{\mathrm{dB}}\approx 20\log_{10}|K|
\]
Si hay un cero en el origen:
\[
|G(j\omega)|_{\mathrm{dB}}\approx 20\log_{10}|K|+20\log_{10}(\omega)
\]
Si hay un polo en el origen:
\[
|G(j\omega)|_{\mathrm{dB}}\approx 20\log_{10}|K|-20\log_{10}(\omega)
\]
Más general:
\[
|G(j\omega)|_{\mathrm{dB}}\approx 20\log_{10}|K|+20n\log_{10}(\omega)
\]
donde $n>0$ indica ceros en el origen y $n<0$ polos en el origen.

% ============================
\section*{11. Fase aproximada por asíntotas}

Reglas básicas:
\[
\text{cero} \Rightarrow +90^\circ \qquad \text{polo} \Rightarrow -90^\circ
\]
\[
s\ \text{arriba} \Rightarrow +90^\circ \qquad s\ \text{abajo} \Rightarrow -90^\circ
\]
Ejemplo con $G(s)=\dfrac{20s(s+10)}{(s+2)(s+25)(s+40)}$:

Contribuciones: $s\!:+90^\circ$,\ $(s+10)\!:+90^\circ$,\ $(s+2)\!:-90^\circ$,\ $(s+25)\!:-90^\circ$,\ $(s+40)\!:-90^\circ$.

Fase por intervalos:
\[
90^\circ,\ 0^\circ,\ 90^\circ,\ 0^\circ,\ -90^\circ
\]

% ============================
\section*{12. Fase exacta}

Para un factor de cero $\left(1+\dfrac{j\omega}{a}\right)$:
\[
+\tan^{-1}\!\left(\frac{\omega}{a}\right)
\]
Para un factor de polo $\dfrac{1}{1+j\omega/a}$:
\[
-\tan^{-1}\!\left(\frac{\omega}{a}\right)
\]
Para cero en el origen $j\omega$: $\angle(j\omega)=+90^\circ$.

Para polo en el origen $\dfrac{1}{j\omega}$: $\angle\!\left(\dfrac{1}{j\omega}\right)=-90^\circ$.

% ============================
\section*{13. Respuesta senoidal en estado estable}

Si la entrada es $x(t)=A\cos(\omega t+\phi)$, la salida es:
\[
y(t)=A|G(j\omega)|\cos\!\left(\omega t+\phi+\angle G(j\omega)\right)
\]
Reglas:
\[
\text{amplitud de salida}=A|G(j\omega)|
\]
\[
\text{fase de salida}=\phi+\angle G(j\omega)
\]
La frecuencia no cambia.

% ============================
\section*{14. Convertir seno a coseno}

Identidades importantes:
\[
\sin(\theta)=\cos(\theta-90^\circ)
\]
\[
\cos(\theta)=\sin(\theta+90^\circ)
\]
\[
\cos(\theta-180^\circ)=-\cos(\theta)
\]
\[
\sin(\theta-90^\circ)=-\cos(\theta)
\]
Ejemplo:
\[
4\sin(\omega t)=4\cos(\omega t-90^\circ)
\]

% ============================
\section*{15. Combinar coseno y seno}

Si tienes $x(t)=a\cos(\omega t)+b\sin(\omega t)$, puedes escribirlo como:
\[
x(t)=R\cos(\omega t-\alpha)
\]
donde:
\[
R=\sqrt{a^2+b^2}, \qquad \alpha=\tan^{-1}\!\left(\frac{b}{a}\right)
\]
porque $R\cos(\omega t-\alpha)=R\cos(\omega t)\cos\alpha+R\sin(\omega t)\sin\alpha$,

entonces $a=R\cos\alpha$ y $b=R\sin\alpha$.

Ejemplo: $3\cos(5t)+4\sin(5t)$
\[
R=\sqrt{3^2+4^2}=5, \qquad \alpha=\tan^{-1}\!\left(\frac{4}{3}\right)=53.13^\circ
\]
\[
x(t)=5\cos(5t-53.13^\circ)
\]

% ============================
\section*{16. Checklist para resolver}

\begin{enumerate}
\item Sustituir $s=j\omega$.
\item Identificar polos, ceros y ganancia $K$.
\item Normalizar factores: $s+a=a(1+s/a)$.
\item Encontrar frecuencias de quiebre.
\item Para magnitud:
\[
|G(j\omega)|=\frac{\text{magnitudes de arriba}}{\text{magnitudes de abajo}}
\]
\item Para fase:
\[
\angle G = \text{fases de arriba} - \text{fases de abajo}
\]
\item Para Bode: sumar $+20$ por cero y $-20$ por polo.
\item Para salida senoidal: multiplicar amplitud y sumar fase.
\end{enumerate}

\end{multicols}

\end{document}
